\section{Methods I: building the dataset}
\label{sec:methods-data}

\subsection{Design principles}

Three principles governed every step, and they explain several choices that would otherwise look unusual.

\begin{enumerate}
  \item \textbf{Declare before answering.} The extraction schema was frozen before extraction started; the
        framework design, the success thresholds, the sensitivity analyses and every reading rule for the
        post-hoc analyses were committed to the repository before the first result they govern was produced. Where
        a rule was added after a result was known, the report says so.
  \item \textbf{The campaign is the unit.} Two papers that report the same participants are one observation of
        the underlying physiology, not two. Every validation in the project holds out whole campaigns, and every
        interval is computed by resampling campaigns.
  \item \textbf{Nothing is typed twice.} The dataset is frozen under a version tag and a SHA-256 hash; the
        framework refuses to run against a file that does not match; every number and figure is regenerated from
        that file by one command (\cref{sec:reproducibility}).
\end{enumerate}

The work was carried out by two leads in parallel tracks that met at a frozen interface: the scientific lead ran
the search, screening and physiological interpretation; the AI framework lead designed the schema, the framework
and the models. The phases and their dates are summarised in \cref{tab:phases}.

\begin{table}[htbp]
\centering
\caption{Project phases. The written report and the presentation slides are phase P5.}
\label{tab:phases}
\small
\begin{tabularx}{\linewidth}{@{}l l Y@{}}
\toprule
Phase & Dates (2026) & Content and outcome \\
\midrule
P0 Kickoff & 4 Sep & Repository, branch model, frozen extraction schema, decisions on the pre-search corpus \\
P1 Literature & 4--11 Sep & Search of four databases, title/abstract and full-text screening, extraction, campaign map \\
P2 Integration & 12--14 Sep & Source verification, reconciliation, quality control; \code{dataset\_v1.0} frozen (14 Sep) \\
P3 Design & 15--18 Sep & Framework design document, configuration, modules, leakage guard, baseline \\
P4 Modelling & 16--21 Sep & Tier 2 (18 Sep); tier 1 and \code{dataset\_v1.1} (19 Sep); time-axis correction, tier 3 with ablations and validation, figures (21 Sep) \\
P5 Report and slides & 22--28 Sep & This document; the presentation deck \\
\bottomrule
\end{tabularx}
\end{table}

\subsection{Literature search}
\label{sec:search}

Four sources were searched on 4 September 2026: PubMed, Scopus, Web of Science Core Collection and the NASA
Technical Reports Server (NTRS). Each query combined two concept blocks, both required: \emph{A}, the unloading
model (bed rest, head-down tilt, dry immersion, limb suspension, simulated microgravity, disuse and their
variants), and \emph{B}, the muscle outcome (atrophy, muscle volume, mass, size, cross-sectional area, lean mass,
thickness, and the names of the principal lower-limb muscles). Imaging-modality and countermeasure blocks were
deliberately left out of the first pass so as not to lose papers whose abstracts omit them. Rodent studies were
excluded in the query where the syntax allowed. Every search was limited to publications from 2013 onwards; older
work enters the dataset only through the pre-search corpus (\cref{sec:corpus}). The NTRS search used five short
keyword queries, because a technical memorandum is not expected to be indexed by the journal databases. The
queries are reproduced verbatim in \cref{app:search}.

Embase, Cochrane CENTRAL, Europe PMC full-text search, trial registries, citation chasing and hand-searching were
not run. Embase was unavailable for lack of institutional access; its absence is a known limitation because it
indexes conference abstracts the other databases do not.

\paragraph{Known-item test} Before a query is trusted, it should return the papers already known to be relevant.
Because of the 2013 date limit, only two of the ten known modelling papers fell inside the search window; both were
retrieved by all three journal databases. The test with the date limit removed was not run.

\subsection{Eligibility criteria}

The criteria in \cref{tab:eligibility} were written before screening began. Two judgement calls were settled in
advance: total limb volume by anthropometry or plethysmography is not a muscle outcome, because fluid shifts
confound it; and cast immobilisation and step reduction are not unloading models for this dataset, because their
mechanical and postural picture differs from bed rest.

\begin{table}[htbp]
\centering
\caption{Eligibility criteria. A study was included when every row held.}
\label{tab:eligibility}
\small
\begin{tabularx}{\linewidth}{@{}l Y@{}}
\toprule
Axis & Requirement \\
\midrule
Population & Humans, adults ($\geq$18 years), any sex, healthy or without a disease that itself causes muscle wasting \\
Exposure & A ground-based unloading model with a stated duration: horizontal or head-down bed rest, dry immersion,
           unilateral lower-limb suspension. Spaceflight eligible but flagged separately \\
Duration & At least 5 days of continuous unloading \\
Outcome & At least one lower-limb muscle-tissue measure at baseline and at one or more later time points: volume,
          cross-sectional area, thickness, or lean or muscle mass, by MRI, CT, DXA, pQCT or ultrasound \\
Reporting & Enough to compute a percentage change per arm: baseline and follow-up values, or a printed change \\
Type & Primary study; reviews were mined for references and never extracted \\
Language & English or German; other languages recorded rather than silently dropped \\
\midrule
Exclusion codes & \code{not\_human}, \code{not\_unloading\_model}, \code{duration\_too\_short},
  \code{no\_muscle\_outcome}, \code{upper\_limb\_only}, \code{no\_baseline\_or\_followup},
  \code{disease\_causes\_wasting}, \code{review\_or\_editorial}, \code{duplicate\_report}, \code{no\_full\_text},
  \code{language} \\
\bottomrule
\end{tabularx}
\end{table}

\subsection{Screening}

Records from the four exports were merged and de-duplicated on DOI and normalised title
(\file{framework/merge\_search\_exports.py}); duplicates were marked rather than deleted. Title and abstract
screening ran in two stages: a deterministic triage against the criteria (\file{framework/triage\_screening.py}),
which sorted records into exclusions with a reason, a priority set and a \emph{maybe} set, followed by a human
reading of the priority set and the top of the \emph{maybe} set. Full texts were retrieved preferentially as
machine-readable XML from Europe PMC, which preserves tables, and otherwise as PDF. The remainder of the
\emph{maybe} set was not screened in the time available; those records are reported as unscreened, not as
excluded.

\subsection{The pre-search corpus}
\label{sec:corpus}

Fifteen documents were held by the team before the search began. Each was read in full and classified with a
written reason: nine were modelling studies, one a methods paper that re-analyses a cohort already extracted, four
were vascular studies retained only as context, and one---a study of boys with Duchenne muscular dystrophy---was
excluded because the disease itself causes wasting. The nine modelling studies are the dataset's only source of
data published before 2013 and include its longest unloading duration, 119 days~\citep{leblanc1992}. They are a
convenience corpus, not a search result, and are counted on their own line in the flow diagram.

\subsection{Extraction schema}
\label{sec:schema}

The schema (\file{data/schema.md}) defines one row as \emph{one comparison against baseline}: one study, one arm,
one muscle, one measurement occasion, carrying the baseline value, the follow-up value and the percentage change.
There are no baseline-only rows. The 61 columns fall into groups for study identity and provenance, unloading
design, arm and participants, what was measured, values, and quality control; muscles are named from a controlled
vocabulary that can be extended only by a change to the schema. Ten rules accompany the schema; the ones that
carry the analysis are:

\begin{itemize}
  \item \textbf{Never impute silently}: a missing value is recorded as missing, and that is itself a fact about
        the literature.
  \item \textbf{Sign convention}: \code{pct\_change} is negative for atrophy, everywhere.
  \item \textbf{Recompute every percentage} where both values are printed, and compare it with the printed change.
  \item \textbf{Composites are not sums}: \code{triceps\_surae} is not the sum of soleus and gastrocnemius, and a
        composite and its components must never enter a model together unresolved.
  \item \textbf{One \code{cohort\_id} per campaign}, not per paper.
  \item \textbf{Recovery is a different question}: a recovery row never enters a duration--response model without
        an explicit flag.
  \item \textbf{\code{n\_analysed}, not \code{n\_arm}, is the weight}, because dropouts are routine.
  \item \textbf{Every row is traceable in under a minute} through \code{source\_file} and \code{page\_ref}.
  \item \textbf{Digitised figures are labelled} as such (\code{data\_source = figure\_digitized}).
\end{itemize}

Twelve validation checks run as a script (\file{framework/validate\_extraction.py}) so that a malformed table fails
before it is committed. The schema is summarised in \cref{app:schema}.

\paragraph{Two estimators in one column} Where a paper prints the mean of the participants' individual percentage
changes and it differs from the value recomputed from the group means by more than rounding, the printed value is
kept (\code{pct\_of\_individual\_means}); otherwise the recomputed value is kept
(\code{pct\_recomputed\_from\_group\_means}). The two are different estimators when the change correlates with
baseline size, and every row records which it carries.

\subsection{Source verification and reconciliation}
\label{sec:verification}

On 4 September 2026 the scientific lead re-extracted every provenance record then present from the original
papers, with AI assistance, and compared the result with the extraction tables. The check was staged: three
studies (26 rows) to test the method, then 113 rows (all 14 figure-derived rows and 99 table-derived rows), then
all 580 records. The outcome is summarised in \cref{tab:verification}.

\begin{table}[htbp]
\centering
\caption{Source verification of the extraction (4 September 2026).}
\label{tab:verification}
\small
\begin{tabular}{@{}lr@{}}
\toprule
Measure & Result \\
\midrule
Provenance records checked & 580 (537 main, 29 partial, 14 figure) \\
Recovered from the source & 580 (575 high confidence, 5 medium) \\
Agree without correction & 556 \\
Flagged & 24 rows, plus 3 duplicated observations \\
Unresolved numeric conflicts & 0 \\
Percent-change gap: median / mean / maximum & 0 / 0.012 / 1.04\pp \\
\bottomrule
\end{tabular}
\end{table}

The flagged rows were resolved on 14 September and each resolution is logged in \file{data/reconciliation\_log.md}.
Eight rows had been coded as MRI but measured lumbar multifidus by ultrasound and were corrected. Sixteen rows
printed a whole-number percentage that differed from the recomputed value; because a whole number disagrees with a
precise value only once the gap exceeds its own rounding, the printed value was kept where the gap exceeded
0.5\pp{} (five rows) and the recomputed value otherwise. Duplicated observations were flagged and counted once, and
two corrections were made to the campaign map.

This check is \emph{AI-assisted independent source verification}. It is not a human double extraction, and because
the assistant had seen some rows earlier it is not fully blinded; accordingly, \code{double\_extracted} remains
\code{FALSE} on every row. The 165 rows from the pre-search corpus were added after the check and have not been
verified by a second person.

\subsection{The campaign map}
\label{sec:cohort-map}

\file{data/cohorts.csv} assigns every study to the campaign that produced its participants, using trial registry
identifiers where they exist and the text of the papers otherwise. Each assignment carries a note and a
\code{verified} flag. Examples of the dependence it captures: the 90-day Long-Term Bed Rest study at MEDES is
reported by four papers in the dataset~\citep{alkner2004,belavy2017,trappe2023,rittweger2013}; the WISE-2005
women's campaign by five~\citep{trappe2007,rogers2025,trappe2023,holt2015,holt2016}; and the McGill 14-day
campaign in older adults by three that share one registry
identifier~\citep{dulac2024,hajjboutros2023,lagace2026}. The map also separates campaigns that an author-based
grouping would merge: the two Berlin bed-rest studies, BBR1 (56 days, horizontal) and BBR2-2 (60 days, head-down),
are distinct campaigns with distinct participants~\citep{belavy2009,miokovic2012}. The resulting 36 campaigns are
listed in \cref{app:campaigns}.

\subsection{Open NASA data}
\label{sec:nasa}

On 19 September 2026, seven campaign folders of derived, individual-participant bed-rest data were downloaded from
the NASA Life Sciences Portal (NLSP) as open US government data, with every file's identifier and SHA-256 recorded
in a manifest (\file{framework/fetch\_nasa\_nlsp.py}). Only one folder was pooled into the dataset. Four of the
seven belong to campaigns the dataset already holds as literature rows---three to the 70-day NASA SPRINT campaign and
one to WISE-2005---and pooling them would have counted the same participants twice. Two carried no phase labels, so
a baseline would have had to be guessed. The seventh, whole-body DXA from the NASA Flight Analog Project's UTMB
Campaign~3~\citep{nlsp2026}, was pooled for the three blocks whose own scan dates confirm a 90-day, 6$^\circ$
head-down bed rest (13 participants). Leg lean mass was aggregated to five cohort-level rows
(\file{framework/extract\_nasa.py}): four occasions during bed rest and one in recovery. These are the only rows
in the dataset computed from individual measurements rather than read from a publication. Three of the four
in-bed-rest occasions (days 30, 45 and 58) rest on a single participant and are weighted accordingly
(\code{n\_analysed = 1}).

\subsection{Build, quality control and freeze}
\label{sec:freeze}

\file{framework/build\_dataset.py} stacks the four extraction tables (main, partial, figure-derived and NASA) with a
\code{source\_table} column in front of the 61 schema columns, leaves out ten rows that are second copies of an
observation another row already carries, and checks the result across tables: unique \code{row\_id}, known
campaigns, \code{pct\_change} within range, and the minimum viable dataset declared in the plan (60 rows, 8
campaigns, 4 muscle groups, 4 durations). Nothing else is filtered: recovery, spaceflight, dry immersion, DXA,
partial and figure-derived rows are all in the file, each marked by its own column, so that the modelling subset
is chosen in the configuration rather than built into the data.

Version 1.0 was frozen and tagged on 14 September 2026 (737 rows, 51 studies, 35 campaigns). Version 1.1, frozen
on 19 September, added the five NASA rows and one campaign, and is the version every result in this report was
fitted against. The build reproduces each version byte for byte from its tables and refuses to overwrite a frozen
file; the loader verifies the SHA-256 on every run. The one schema change v1.1 required---a \code{repository}
value for \code{data\_source}, so that a row computed from an open archive can say so---is recorded in the
reconciliation log. \draftnote{this schema change carries one lead's signature; the frozen schema requires both.}

\subsection{The modelling subset}
\label{sec:subset}

The modelling subset is defined by three predicates declared in \file{framework/config.yaml} and applied in order:

\begin{enumerate}
  \item \textbf{During unloading only} (\code{phase = bed\_rest}). Recovery rows answer a different question and
        are reserved for a reconditioning sensitivity analysis.
  \item \textbf{Lower-limb muscles only.} Multifidus, lumbar erector spinae, quadratus lumborum, psoas and
        iliopsoas are dropped: trunk muscles unload differently and would answer another question inside the same
        coefficient.
  \item \textbf{One tissue per measurement occasion.} A measurement occasion is a campaign, arm, scan day,
        modality and outcome. Where a paper reports both a composite (for example the quadriceps) and components
        that cover it (its four heads), the components are kept and the composite dropped, because a model given
        both fits the same tissue twice. Where no components are present the composite is kept and flagged.
        Whole-segment measurements (whole thigh, whole calf, whole lower limb) are never decomposed or suppressed.
\end{enumerate}

After predicates 1 and 2 the subset holds 425 rows from 42 studies and 32 campaigns, 299 from control arms and 126
from countermeasure arms, with 13 planned durations and scans on 24 distinct days between day 5 and day 119. Of
these rows, 369 are MRI (314 MRI volume), 29 DXA, 15 CT and 12 ultrasound. Both a composite and its components were
present on 20 of the 93 measurement occasions.

\paragraph{Two subsets, for two claims} Seven campaigns report only whole-segment measurements, and six of them
sit between 5 and 21 days, which is where the duration curve bends most. Excluding whole-segment rows would cost a
fifth of the campaigns and the short-duration evidence with them. Subset A, used for the duration--response and for
every tier-2 and tier-3 comparison, therefore keeps them as a \code{whole\_limb} level of the muscle family (346
rows, 32 campaigns); subset B, used for the muscle ranking, keeps named muscles only (304 rows, 25 campaigns).
\Cref{tab:subsets} in \cref{sec:results-data} walks the dataset down to both.

\paragraph{The sample size is 32, not 346} The largest campaign, MEDES LTBR, contributes 86 rows to subset A and
the two Berlin campaigns another 82 between them. Every design decision in the framework follows from treating
the number of independent campaigns as the effective sample size.

\subsection{Target and exposure}
\label{sec:target}

\paragraph{Target} The target is \code{pct\_change}, the percentage change in muscle size from the same arm's
baseline, negative for loss. In subset A it ranges from $-30.0\%$ to $+24.3\%$ (mean $-8.0\%$, median $-6.8\%$);
the positive values are real, coming from countermeasure arms and small muscles under compensatory loading. No
transformation is applied, so that coefficients read directly in percentage points. The dispersion reported with
each value is heterogeneous: of the 425 lower-limb unloading rows, 290 carry a standard deviation, 88 a standard
error, 6 a 95\% confidence interval, 2 an interquartile range and 39 nothing; and most papers report the spread
of the baseline rather than of the change.

\paragraph{Exposure: the day of the scan} The exposure is the number of days of unloading at the time of the scan
(\code{timepoint\_days}), not the campaign's planned length (\code{duration\_days}). The first runs of tiers 1 and
2 used the planned length, which placed 200 of the 346 rows in subset A---every scan taken before bed rest ended,
120 of them more than a week before---at the end of their campaign, so that a day-14 scan in a 56-day campaign
entered the model as 56 days. The error was found and corrected on 21 September 2026; it is a correction of the
implementation, since the design already named the scan day as a within-campaign variable, not a new analysis
choice. The exposure is now declared once in the configuration and read through one function by every tier. The
effect of the correction is reported in \cref{sec:results-correction}.

\subsection{Muscle classification}
\label{sec:muscle-map}

The schema's \code{muscle\_function\_class} was empty in every extracted row, and \code{composite\_of} was populated
on only 50 of the 189 composite rows in the modelling subset. Both are therefore supplied by a reviewable table,
\file{data/muscle\_map.csv}, joined at load time rather than written into the frozen dataset, so that an assignment
can be argued with and re-run without a new dataset version. It gives each of the 51 muscles a family (plantar
flexors, dorsiflexors, evertors, knee extensors, knee flexors, hip extensors, flexors, abductors, adductors and
rotators, trunk extensors, whole segment), a functional class, its components where it is a composite, a
fibre-type profile and a written rationale; the loader refuses to run if the dataset names a muscle the table does
not cover. The full table is in \cref{app:musclemap}.

Function was assigned by mechanical role first and fibre-type composition second~\citep{johnson1973}: 19 muscles
are antigravity extensors, 8 flexors and 24 of mixed function. Three judgement calls are flagged for review.
Rectus femoris is \emph{mixed}, not an antigravity extensor, because it extends the knee but flexes the hip and
loses less than the monoarticular vasti. Tibialis anterior is a \emph{flexor} despite being slow-fibre dominant,
because it lifts the foot rather than carrying body weight; it is the cleanest test in the corpus of role against
fibre type. And the deep posterior compartment---flexor digitorum longus, flexor hallucis longus and tibialis
posterior---sits with the antigravity extensors, because what these muscles do at the ankle is plantar flexion.
\draftnote{the classification awaits the scientific lead's physiological sign-off, which the plan makes a blocking
condition for quoting coefficients by class.}
